\begin{statementbox}
\textbf{\textcolor{CStatement}{条件}}
椭圆 $x^2/a^2+y^2/b^2=1$，$a>b>0$，$c=\sqrt{a^2-b^2}$；
焦点为 $F_1(-c,0)$、$F_2(c,0)$。直线经过其中一个焦点，与椭圆交于 $A,B$。
记倾斜角 $\theta\in[0,\pi)$（相对 $x$ 轴正方向），弦长 $L=AB$。

\medskip
\textbf{\textcolor{CStatement}{焦点弦长}}
\[
 \boxed{L=\frac{2ab^2}{a^2-c^2\cos^2\theta}
 =\frac{2h}{1-e^2\cos^2\theta}},\qquad e=\frac ca,\quad h=\frac{b^2}{a}.
\]
$a$ 是半长轴，$c$ 是半焦距，$e$ 是离心率，$h$ 是半通径。
左、右焦点使用同一弦长公式，且分母至少为 $b^2>0$。

\medskip
\textbf{\textcolor{CStatement}{两侧焦半径}}
取右焦点 $F_2$，将 $A$ 选在方向 $(\cos\theta,\sin\theta)$ 一侧，
$B$ 在相反一侧。令 $r=F_2A$、$s=F_2B$，则
\[
 r=\frac{b^2}{a+c\cos\theta},\qquad
 s=\frac{b^2}{a-c\cos\theta},\qquad \boxed{L=r+s}.
\]
两分母均为正（$a>c$），两段通常不等，垂直长轴时相等。

\medskip
\textbf{\textcolor{CStatement}{通径及最值}}
\[
 \boxed{\frac{2b^2}{a}\le L\le2a}.
\]
最小值在 $\theta=90^\circ$ 时取得，是过焦点且垂直长轴的通径；
最大值在 $\theta=0^\circ$ 时取得，是长轴弦。

\medskip
\textbf{\textcolor{CStatement}{斜率与长轴方向}}
有限斜率 $k$ 满足 $\cos^2\theta=1/(1+k^2)$，故
\[
 L=\frac{2ab^2(1+k^2)}{a^2(1+k^2)-c^2}.
\]
竖直线直接用通径。长轴在 $y$ 轴时，改用弦与长轴的角 $\varphi$：
\[
 L=\frac{2ab^2}{a^2-c^2\cos^2\varphi}.
\]
其中 $a$ 始终是半长轴；若仍给相对 $x$ 轴的倾斜角 $\theta$，则 $\cos^2\varphi=\sin^2\theta$。
\end{statementbox}
